\begin{abstract}
Random Forests are widely used for hybrid Wi-Fi round-trip time (RTT) and received signal strength (RSS) fingerprinting, but prior work considers all features at every split ($\rho=1$), which reduces the forest to bagged regression trees. This paper studies the fraction $\rho$ of features considered per split along three questions: how it depends on radio-map density, how it should be selected, and how uncertainty regions should be calibrated when each reference point (RP) contributes many near-identical scans. On three public environments, $\rho=0.7$ reduces the error of the full-feature forest by 14.4\%, 9.2\%, and 5.2\% on RP-disjoint test sets, in each of ten seeds. \emph{Radio-map density:} $\rho=1$ is never optimal; as the training map is thinned, the optimal ratio falls from 0.5--0.85 to 0.2--0.35 and the gain over $\rho=1$ reaches 26\%, whereas extremely randomized trees do not benefit from subsampling. \emph{Ratio selection:} out-of-bag error and cross-validation favor ratios that are too small, whereas the fixed ratio $\rho=0.7$ stays within 4.1\% of the best ratio in every environment ($\rho=1$: 19.7\%). It also outperforms $\rho=1$ in 74 of 80 comparisons on an independent two-day dataset, once a mirrored coordinate frame between its days is corrected. \emph{Calibration:} with 21 calibration RPs, scan-level split-conformal regions under-cover on average, whereas group calibration with one scan per RP attains nominal coverage with 22--30\% larger regions and needs at least $1/\alpha-1$ calibration RPs. Feature subsampling should thus match survey density, and validation and calibration should treat RPs, not scans, as the unit.
\end{abstract}

\begin{keyword}
indoor positioning \sep Wi-Fi round-trip time \sep fingerprinting \sep random forest \sep feature subsampling \sep radio-map density \sep conformal prediction
\end{keyword}
